%!TeX root = ../main.tex
\chapter{Introdução}

% Um emparelhamento em um grafo não-dirigido é definido como um conjunto de arestas sem pontas em comum. Em outras palavras, um emparelhamento é um conjunto $M$ de arestas que satisfaz a seguinte propriedade: o grau de cada vértice no subconjunto $M$ é no máximo 1 \cite{imeUSP_matching}. Os principais problemas relacionados a emparelhamentos geralmente envolvem grafos bipartidos, aqueles que podem ser particionados em dois conjuntos independentes $U$ e $V$, onde todas as arestas conectam vértices de $U$ com vértices de $V$.

% Os problemas de emparelhamento têm aplicações significativas em várias áreas, como biologia computacional, onde podem ser empregados na diferenciação e classificação de estruturas proteicas \cite{Taylor2002}; redes sociais, para identificação de comunidades ou grupos \cite{wenfei_socialnetwork}; e visão computacional, para correspondências entre elementos de conjuntos distintos, como pontos em imagens ou vértices em malhas tridimensionais \cite{haller2022comparativestudygraphmatching}.

% Resolver problemas de emparelhamento em grafos é um desafio devido à sua complexidade computacional, frequentemente NP-completa, exigindo métodos inovadores para equilibrar precisão e eficiência. Estratégias clássicas, como o Problema de Atribuição Quadrática(QAP), modelam o emparelhamento como uma questão de otimização combinatória, empregando relaxações espectrais e heurísticas para encontrar soluções aproximadas de forma eficiente \cite{ijcai2020p694}. Por outro lado, avanços em transporte ótimo duplamente estocástico, como o algoritmo GOAT, têm demonstrado melhorias em robustez e velocidade, especialmente em grafos maiores \cite{saadeldin_goat}.

% Estudos como esses ilustram a importância de investigar e comparar métodos computacionais que abordem a complexidade inerente ao problema de emparelhamento em grafos, contribuindo para aplicações como aprendizado de máquina e visão computacional. Essa abordagem torna-se essencial para desenvolver soluções adaptadas a diferentes cenários e com impacto direto na melhoria de algoritmos existentes.

% Diante desse cenário, este trabalho visa sistematizar os principais problemas e métodos computacionais associados ao emparelhamento em grafos. Ele busca oferecer uma comparação detalhada entre suas aplicações e limitações, contribuindo para o entendimento aprofundado das ferramentas disponíveis e identificando oportunidades para o desenvolvimento de novas soluções. Essa abordagem torna-se ainda mais relevante à luz da importância prática desses algoritmos para resolver problemas reais em grande escala, conforme evidenciado por estudos recentes \cite{MIT2021}.

Na teoria dos grafos e otimização combinatória, um emparelhamento em um grafo não-dirigido $G=(V, E)$ é definido formalmente como um subconjunto de arestas $M \subseteq E$ tal que nenhuma aresta em $M$ compartilhe um vértice comum. Em termos precisos, para todo vértice $v \in V$, o grau de $v$ no subgrafo induzido por $M$ é no máximo 1. Esta definição fundamental é apresentada por \citeonline{lovasz-cap1} em sua obra seminal sobre a teoria dos emparelhamentos, onde também se estabelece a distinção clássica entre emparelhamentos em grafos bipartidos — cujos vértices podem ser divididos em dois conjuntos disjuntos $U$ e $V$ — e em grafos gerais.

Os problemas de emparelhamento têm aplicações significativas em várias áreas, como biologia computacional, onde podem ser empregados na diferenciação e classificação de estruturas proteicas \cite{Taylor2002}; redes sociais, para identificação de comunidades ou grupos \cite{wenfei_socialnetwork}; e visão computacional, para correspondências entre elementos de conjuntos distintos, como pontos em imagens ou vértices em malhas tridimensionais \cite{haller2022comparativestudygraphmatching}.

É crucial, no entanto, fazer uma distinção precisa quanto à complexidade computacional desses problemas. Ao contrário do que uma análise superficial poderia sugerir, o problema de encontrar um emparelhamento de cardinalidade ou peso máximo pertence à classe de complexidade \textbf{P} (tempo polinomial), sendo solúvel por algoritmos clássicos como o Método Húngaro ou o Algoritmo de \textit{Blossom} de Edmonds \cite{schrijver}. O verdadeiro desafio computacional, frequentemente classificado como \textbf{NP-Difícil}, surge em variações estruturais mais complexas, como o Emparelhamento Tridimensional ou o Problema da Atribuição Quadrática (QAP). Segundo \citeonline{garey-jhonson} e \citeonline{burkard}, o QAP modela situações onde o custo depende da interação entre os pares escolhidos, tornando métodos exatos inviáveis para grandes instâncias e exigindo o uso de relaxações ou heurísticas avançadas.

Diante desse cenário, este trabalho visa sistematizar os principais problemas e métodos computacionais associados ao emparelhamento em grafos. Ele busca oferecer uma comparação detalhada entre suas aplicações e limitações, contribuindo para o entendimento aprofundado dos algoritmos disponíveis e identificando oportunidades para o desenvolvimento de novas soluções. Essa abordagem torna-se ainda mais relevante à luz da importância prática desses algoritmos para resolver problemas reais em grande escala, conforme evidenciado por estudos recentes \cite{MIT2021}.

\section{Objetivos}

Os objetivos deste trabalho dividem-se em um objetivo geral, que expressa o propósito do estudo como um todo, e em objetivos específicos, que o decompõem nas etapas necessárias para alcançá-lo.

\subsection{Objetivo geral}

Este trabalho tem como objetivo principal levantar e organizar os principais problemas de emparelhamento em grafos e os métodos computacionais presentes na literatura para resolvê-los, implementando esses métodos e comparando-os experimentalmente quanto ao desempenho e à qualidade das soluções.

\subsection{Objetivos específicos}

\begin{itemize}
    \item Mapear as variações e características dos principais problemas de emparelhamento em grafos presentes na literatura.
    \item Classificar esses problemas e os métodos computacionais que os resolvem, segundo critérios como a estrutura do grafo, o objetivo de otimização e o tipo de método (exato, aproximado ou heurístico).
    \item Implementar os algoritmos selecionados, adaptando implementações públicas ou, quando não houver implementação adequada, partindo da descrição bibliográfica.
    \item Criar um ambiente de testes unificado e, nele, comparar os algoritmos quanto ao tempo de execução, à complexidade assintótica de tempo e espaço e à qualidade das soluções encontradas.
    \item Organizar os resultados em um formato que facilite o acesso e a compreensão dos tópicos abordados por pesquisadores e profissionais.
\end{itemize}

\section{Delimitação do estudo}

Os problemas de emparelhamento em grafos podem ser classificados de diversas maneiras, como sendo em grafos bipartidos ou não-bipartidos, ponderados ou não-ponderados, direcionados ou não-direcionados, problemas em multigrafos, problemas em hipergrafos, entre outros. Este trabalho focará em problemas que envolvem grafos menos complexos, excluindo multigrafos e hipergrafos. Também ficam fora do escopo outras variações do emparelhamento em grafos gerais, como o b-emparelhamento, o emparelhamento em matroides e o problema dos colegas de quarto estáveis (\textit{Stable Roommates Problem}).

\section{Organização do trabalho}

O restante deste trabalho está organizado da seguinte forma:

\begin{itemize}
    \item \textbf{Capítulo~\ref{cap:fundamentacao}:} apresenta a fundamentação teórica: definições sobre grafos, algoritmos de busca e de caminhos mais curtos, redes de fluxo, o problema geral do emparelhamento, programação linear e dualidade, e a classificação dos métodos em exatos, aproximados e heurísticos.
    \item \textbf{Capítulo~\ref{cap:bipartidos}:} descreve os problemas de emparelhamento em grafos bipartidos (cardinalidade máxima, atribuição, atribuição gargalo e emparelhamento estável), com suas propriedades, seus algoritmos em pseudocódigo e uma tabela-resumo por problema.
    \item \textbf{Capítulo~\ref{cap:gerais}:} faz o mesmo para os problemas em grafos gerais (cardinalidade máxima, peso máximo e emparelhamento induzido máximo).
    \item \textbf{Capítulo~\ref{cap:metodologia}:} descreve o ambiente de testes construído para comparar os algoritmos: arquitetura, interface padrão, geradores de instâncias, protocolo de medição, critérios de comparação e origem de cada implementação.
    \item \textbf{Capítulo~\ref{cap:resultados}:} apresenta e discute os resultados obtidos no ambiente de testes, quanto ao tempo de execução, à qualidade das soluções e ao crescimento empírico frente à complexidade teórica.
    \item \textbf{Capítulo~\ref{cap:trabalhos_futuros}:} aponta as extensões deste trabalho que ficaram como trabalhos futuros.
    \item \textbf{Apêndice~\ref{ap:compilados}:} reúne as tabelas de tempo das implementações de núcleo compilado, usadas como referência.
\end{itemize}
